% resolucoes_exame.tex

\section*{Exercício 1 — Equação Quadrática}
Enunciado: Resolva \(x^2 - 5x + 6 = 0\).



\[
x^2 - 5x + 6 = 0
\]



\textbf{Passo 1:} Identificar coeficientes: \(a=1, b=-5, c=6\).  
\textbf{Passo 2:} Calcular discriminante:


\[
\Delta = b^2 - 4ac = 25 - 24 = 1
\]


\textbf{Passo 3:} Fórmula de Bhaskara:


\[
x = \frac{-b \pm \sqrt{\Delta}}{2a} = \frac{5 \pm 1}{2}
\]


\textbf{Passo 4:} Raízes:


\[
x_1 = 2, \quad x_2 = 3
\]



---

\section*{Exercício 2 — Limite Trigonométrico}
Enunciado: Calcule \(\lim_{x \to 0} \frac{\sin(x)}{x}\).



\[
\lim_{x \to 0} \frac{\sin(x)}{x}
\]



\textbf{Passo 1:} Usar série de Taylor:


\[
\sin(x) = x - \frac{x^3}{3!} + \frac{x^5}{5!} - \dots
\]


\textbf{Passo 2:} Dividir por \(x\):


\[
\frac{\sin(x)}{x} = 1 - \frac{x^2}{3!} + \frac{x^4}{5!} - \dots
\]


\textbf{Passo 3:} Limite:


\[
\lim_{x \to 0} \frac{\sin(x)}{x} = 1
\]



---

\section*{Exercício 3 — Derivada de Polinômio}
Enunciado: Calcule \(\frac{d}{dx}(x^3 - 4x^2 + x - 7)\).



\[
f(x) = x^3 - 4x^2 + x - 7
\]



\textbf{Passo 1:} Regra da potência: \(\frac{d}{dx}(x^n) = nx^{n-1}\).  
\textbf{Passo 2:} Derivar cada termo:


\[
f'(x) = 3x^2 - 8x + 1
\]



---

\section*{Exercício 4 — Integral Definida}
Enunciado: Calcule \(\int_0^1 (2x+1)\,dx\).



\[
\int_0^1 (2x+1)\,dx
\]



\textbf{Passo 1:} Separar termos:


\[
\int_0^1 2x\,dx + \int_0^1 1\,dx
\]


\textbf{Passo 2:} Calcular integrais:


\[
\int 2x\,dx = x^2, \quad \int 1\,dx = x
\]


\textbf{Passo 3:} Avaliar limites:


\[
[x^2]_0^1 + [x]_0^1 = (1-0) + (1-0) = 2
\]



---

\section*{Exercício 5 — Sistema Linear}
Enunciado: Resolva


\[
\begin{cases}
x + y = 5 \\
2x - y = 1
\end{cases}
\]



\textbf{Passo 1:} Da primeira equação: \(y = 5 - x\).  
\textbf{Passo 2:} Substituir na segunda:


\[
2x - (5 - x) = 1 \Rightarrow 3x - 5 = 1
\]


\textbf{Passo 3:} Resolver:


\[
3x = 6 \Rightarrow x = 2
\]


\textbf{Passo 4:} Substituir em \(y = 5 - x\):


\[
y = 3
\]



\textbf{Resultado:} \(x=2, y=3\).
